\subsubsection{边界方程塔、solenoid 包络与相位--寄存双轴预算}\label{subsubsec:conclusion-boundary-padic-solenoid-phase-ledger}
沿用推论 \ref{cor:conclusion-boundary-cycle-rank-external-info-lower-bound} 的记号，记
\[
r:=r(f)=\operatorname{rank}_{\ZZ}H_1(\mathsf B_f;\ZZ)=A(f)-|X|+c(\mathsf B_f).
\]
前述模 \(p\) 仿射纤维、pro-\(p\) 寄存最小性与 solenoid 终对象性可进一步压缩为一条统一结构链：回路秩 \(r\) 同时控制有限层边界方程的解纤维基数、逐层精确选择器的最小离散复杂度、其 \(p\)-进与阿德尔完成、有限素数联合读出的相位/寄存乘法律，以及通用连通包络上的刚性与无穷素数尾的结构性失明。

\begin{theorem}[\(p\)-进边界方程塔与精确秩]\label{thm:conclusion-boundary-padic-tower-exact-rank}
\leanverified{paper\_conclusion\_boundary\_padic\_tower\_exact\_rank}
取 \(\mathsf B_f\) 的任意定向关联矩阵
\[
B\in M_{|X|\times A(f)}(\ZZ).
\]
对任意素数 \(p\)、整数 \(\ell\ge 1\) 与
\[
\delta_\ell\in \operatorname{im}\!\Bigl(B:(\ZZ/p^\ell\ZZ)^{A(f)}\to (\ZZ/p^\ell\ZZ)^{X}\Bigr),
\]
定义高阶边界方程解纤维
\[
\mathcal F_{p^\ell}(\delta_\ell)
:=
\left\{
\varphi\in (\ZZ/p^\ell\ZZ)^{A(f)}:\ B\varphi=\delta_\ell
\right\}.
\]
则 \(\mathcal F_{p^\ell}(\delta_\ell)\) 是一个秩 \(r\) 的仿射 \((\ZZ/p^\ell\ZZ)\)-模；特别地，
\[
\boxed{\;
\card{\mathcal F_{p^\ell}(\delta_\ell)}=p^{\ell r}.
\;}
\]
若 \((\delta_\ell)_{\ell\ge 1}\) 在模约化下相容，则逆极限
\[
\mathcal F_{p^\infty}(\delta)
:=
\varprojlim_{\ell}\mathcal F_{p^\ell}(\delta_\ell)
\]
非空，并且是 \(\ZZ_p^{\,r}\) 上的主齐次空间；选定任一基点后，有非规范同构
\[
\mathcal F_{p^\infty}(\delta)\cong \ZZ_p^{\,r}.
\]
\end{theorem}
\begin{proof}
有限图关联矩阵的 Smith 标准形给出整数可逆矩阵 \(U,V\) 使
\[
UBV=\operatorname{diag}\!\bigl(
\underbrace{1,\dots,1}_{|X|-c(\mathsf B_f)},
\underbrace{0,\dots,0}_{r}
\bigr).
\]
这里非零不变量全为 \(1\)，而零不变量的个数恰为
\(
A(f)-|X|+c(\mathsf B_f)=r
\)。
将此式模 \(p^\ell\) 约化后，\(B\) 在 \(\ZZ/p^\ell\ZZ\) 上与同一对角形等价，因此其核同构于
\[
(\ZZ/p^\ell\ZZ)^{r}.
\]
若 \(\delta_\ell\) 属于像，则 \(\mathcal F_{p^\ell}(\delta_\ell)\) 非空，且为该核的一个平移，从而是秩 \(r\) 的仿射模，并满足
\[
\card{\mathcal F_{p^\ell}(\delta_\ell)}
=
\card{(\ZZ/p^\ell\ZZ)^r}
=
p^{\ell r}.
\]
当 \(\ell=1\) 时，这正退化为推论 \ref{cor:conclusion-boundary-cycle-rank-external-info-lower-bound} 的模 \(p\) 结论。

若 \((\delta_\ell)\) 相容，则在 Smith 坐标下，约化映射
\[
\mathcal F_{p^{\ell+1}}(\delta_{\ell+1})\to \mathcal F_{p^\ell}(\delta_\ell)
\]
逐坐标表现为
\(
(\ZZ/p^{\ell+1}\ZZ)^r\twoheadrightarrow(\ZZ/p^\ell\ZZ)^r
\)
上的平移，故必为满射。
于是逆极限非空，且差分群为
\[
\varprojlim_\ell(\ZZ/p^\ell\ZZ)^r=\ZZ_p^{\,r}.
\]
故 \(\mathcal F_{p^\infty}(\delta)\) 为 \(\ZZ_p^{\,r}\)-torsor。
证毕。
\end{proof}

\leanverified{paper\_conclusion\_boundary\_padic\_tower\_exact\_slope\_rigidity}
\begin{corollary}[\(p\)-进边界方程塔的精确斜率刚性]\label{cor:conclusion-boundary-padic-tower-exact-slope-rigidity}
在定理 \ref{thm:conclusion-boundary-padic-tower-exact-rank} 的设定下，若 \((\delta_\ell)_{\ell\ge 1}\) 在模约化下相容，则对每个 \(\ell\ge 1\) 都有
\[
\frac{\card{\mathcal F_{p^{\ell+1}}(\delta_{\ell+1})}}{\card{\mathcal F_{p^\ell}(\delta_\ell)}}=p^r.
\]
并且约化映射
\[
\pi_{\ell+1,\ell}:\mathcal F_{p^{\ell+1}}(\delta_{\ell+1})\longrightarrow \mathcal F_{p^\ell}(\delta_\ell)
\]
是满射，其每条纤维恰有 \(p^r\) 个点。特别地，
\[
\boxed{\;
r=\frac{1}{\ell}\log_p\card{\mathcal F_{p^\ell}(\delta_\ell)}
\qquad(\ell\ge 1).
\;}
\]
\end{corollary}
\begin{proof}
由定理 \ref{thm:conclusion-boundary-padic-tower-exact-rank}，
\[
\card{\mathcal F_{p^\ell}(\delta_\ell)}=p^{\ell r},
\qquad
\card{\mathcal F_{p^{\ell+1}}(\delta_{\ell+1})}=p^{(\ell+1)r},
\]
故层间比值为 \(p^r\)。同一定理的证明还说明
\(
\pi_{\ell+1,\ell}
\)
在 Smith 坐标下逐坐标表现为
\[
(\ZZ/p^{\ell+1}\ZZ)^r\twoheadrightarrow (\ZZ/p^\ell\ZZ)^r,
\]
因而为满射，且每条纤维基数恰为 \(p^r\)。最后一式由
\(
\card{\mathcal F_{p^\ell}(\delta_\ell)}=p^{\ell r}
\)
直接取 \(\log_p\) 得到。证毕。
\end{proof}

\begin{theorem}[逐层精确选择器的最小离散复杂度]\label{thm:conclusion-boundary-exact-selector-complexity}
固定素数 \(p\) 与层级 \(\ell\ge 1\)。令
\[
B_\ell:(\ZZ/p^\ell\ZZ)^{A(f)}\to (\ZZ/p^\ell\ZZ)^X
\]
为边界矩阵 \(B\) 的模 \(p^\ell\) 约化。称映射
\[
\Sigma_\ell:(\ZZ/p^\ell\ZZ)^{A(f)}\longrightarrow \operatorname{im}(B_\ell)\times R_\ell,
\qquad
\Sigma_\ell(\varphi)=\bigl(B_\ell\varphi,\operatorname{lab}_\ell(\varphi)\bigr)
\]
为一条逐层精确选择器，若 \(\Sigma_\ell\) 为单射。则最小离散寄存器规模满足
\[
\boxed{\;
\min \card{R_\ell}=p^{\ell r}.
\;}
\]
等价地，最小比特预算恰为
\[
\boxed{\;
\mathsf{bits}_{\min}(\ell,p)=\ell r\log_2 p.
\;}
\]
\leanverified{paper\_conclusion\_boundary\_exact\_selector\_complexity}
\end{theorem}
\begin{proof}
任取一条非空纤维 \(\mathcal F_{p^\ell}(\delta_\ell)\)。若 \(\Sigma_\ell\) 为单射，则其在该纤维上的限制仍为单射，故
\[
\card{R_\ell}\ge \card{\mathcal F_{p^\ell}(\delta_\ell)}=p^{\ell r}
\]
由定理 \ref{thm:conclusion-boundary-padic-tower-exact-rank} 得到。

反向取 Smith 标准形坐标
\[
UBV=
\operatorname{diag}\!\bigl(
\underbrace{1,\dots,1}_{|X|-c(\mathsf B_f)},
\underbrace{0,\dots,0}_{r}
\bigr),
\]
并在 \((\ZZ/p^\ell\ZZ)^{A(f)}\) 中写
\[
V^{-1}\varphi=(u,v),
\qquad
u\in(\ZZ/p^\ell\ZZ)^{|X|-c(\mathsf B_f)},
\quad
v\in(\ZZ/p^\ell\ZZ)^r.
\]
于是 \(B_\ell\varphi\) 唯一决定 \(u\)，而末 \(r\) 个坐标 \(v\) 正参数化核方向。取
\[
R_\ell:=(\ZZ/p^\ell\ZZ)^r,
\qquad
\operatorname{lab}_\ell(\varphi):=v.
\]
若 \(\Sigma_\ell(\varphi)=\Sigma_\ell(\psi)\)，则 \(B_\ell\varphi=B_\ell\psi\) 给出两者的 \(u\)-坐标相同，而 \(\operatorname{lab}_\ell(\varphi)=\operatorname{lab}_\ell(\psi)\) 给出 \(v\)-坐标相同，因此 \(\varphi=\psi\)。故上界可达，最小寄存器规模即为 \(p^{\ell r}\)。取二进对数即得最小比特预算。证毕。
\end{proof}

\begin{theorem}[有限素数联合读出的精确选择器乘法律]\label{thm:conclusion-boundary-finite-prime-selector-product-law}
\leanverified{paper\_conclusion\_boundary\_finite\_prime\_selector\_product\_law}
设 \(S\subset \mathbb P\) 为有限素数集，\(\boldsymbol\ell=(\ell_p)_{p\in S}\) 为正整数向量，并记
\[
M_S:=\prod_{p\in S}p^{\ell_p}.
\]
对每个 \(p\in S\) 选定一条非空纤维 \(\mathcal F_{p^{\ell_p}}(\delta_{\ell_p}^{(p)})\)，并定义联合纤维
\[
\mathcal F_{S,\boldsymbol\ell}(\delta)
:=
\prod_{p\in S}\mathcal F_{p^{\ell_p}}(\delta_{\ell_p}^{(p)}).
\]
则区分联合纤维内点所需的最小离散寄存器规模恰为
\[
\boxed{\;
\min \card{R_{S,\boldsymbol\ell}}=M_S^{\,r}.
\;}
\]
等价地，其最小比特预算满足
\[
\boxed{\;
\mathsf{bits}_{\min}(S,\boldsymbol\ell)=r\log_2 M_S.
\;}
\]
\end{theorem}
\begin{proof}
由定理 \ref{thm:conclusion-boundary-padic-tower-exact-rank}，
\[
\card{\mathcal F_{S,\boldsymbol\ell}(\delta)}
=
\prod_{p\in S}\card{\mathcal F_{p^{\ell_p}}(\delta_{\ell_p}^{(p)})}
=
\prod_{p\in S}p^{\ell_p r}
=
M_S^{\,r}.
\]
任何联合精确选择器在该联合纤维上的限制都必须是单射，故其寄存器规模至少为 \(M_S^r\)。

反向对每个 \(p\in S\) 取定理 \ref{thm:conclusion-boundary-exact-selector-complexity} 中的最优寄存器
\(
(\ZZ/p^{\ell_p}\ZZ)^r
\)，并作直积
\[
R_{S,\boldsymbol\ell}
:=
\prod_{p\in S}(\ZZ/p^{\ell_p}\ZZ)^r.
\]
逐素数记录标签即可得到联合精确选择器，其寄存器规模正好为
\[
\prod_{p\in S}p^{\ell_p r}=M_S^r.
\]
故下界可达。取二进对数便得到最小比特预算公式。证毕。
\end{proof}

\begin{theorem}[阿德尔拼接与通用连通包络]\label{thm:conclusion-boundary-adelic-solenoid-envelope}
对每个素数 \(p\)，取一条相容的 \(p\)-进边界方程塔
\(
\mathcal F_{p^\infty}(\delta^{(p)})
\)。
定义阿德尔乘积
\[
\mathcal F_{\widehat{\ZZ}}(\delta)
:=
\prod_{p}\mathcal F_{p^\infty}(\delta^{(p)}).
\]
则 \(\mathcal F_{\widehat{\ZZ}}(\delta)\) 是
\[
\widehat{\ZZ}^{\,r}
=
\prod_{p}\ZZ_p^{\,r}
\]
上的主齐次空间。

进一步，存在典范短正合列
\[
\boxed{\;
0\longrightarrow \widehat{\ZZ}^{\,r}
\longrightarrow \widehat{\QQ}^{\,r}
\xrightarrow{\ \pi_{\QQ,r}\ }
\TT^{r}
\longrightarrow 0,
\;}
\]
并且 \((\widehat{\QQ}^{\,r},\pi_{\QQ,r})\) 在如下范畴中是终对象：对象为满足下列条件的对偶对 \((K,\pi)\)，其中 \(K\) 为紧连通阿贝尔群、\(\pi:K\twoheadrightarrow\TT^r\) 为满射，且其对偶离散群 \(A=\widehat K\) 满足
\begin{enumerate}
  \item \(A\) 无挠；
  \item 对每个素数 \(p\)，乘 \(p\) 映射 \([p]_A\) 为双射；
  \item \(\dim_{\QQ}(A\otimes_{\ZZ}\QQ)=r\)。
\end{enumerate}
等价地，任何满足这些条件的 \(K\) 都与 \(\widehat{\QQ}^{\,r}\) 唯一同构，并与到 \(\TT^r\) 的投影相容。
\end{theorem}
\begin{proof}
第一部分由定理 \ref{thm:conclusion-boundary-padic-tower-exact-rank} 逐素数取直积即得。

对第二部分，取超自然数
\(
\mathfrak N=\prod_{p}p^\infty
\)。
则定义 \ref{def:cdim-highrank-supernatural-localization} 给出的群满足
\(
A_{r,\mathfrak N}=\QQ^r
\) 且
\(
\ZZ_{\mathfrak N}=\widehat{\ZZ}
\)。
因此定理 \ref{thm:cdim-highrank-phase-register-minimality} 在该情形下直接给出
\[
0\to \widehat{\ZZ}^{\,r}\to \widehat{\QQ}^{\,r}\xrightarrow{\pi_{\QQ,r}}\TT^r\to 0.
\]

现设 \((K,\pi)\) 满足定理中的三条条件，记 \(A=\widehat K\)。
由 \([p]_A\) 对每个素数皆为双射，知 \(A\) 为可除群；再与无挠性合并，得到 \(A\) 是一个 \(\QQ\)-向量空间。
又 \(\dim_{\QQ}(A\otimes_{\ZZ}\QQ)=r\)，故 \(A\cong \QQ^r\)。

另一方面，\(\pi:K\twoheadrightarrow\TT^r\) 对偶化后给出一个秩为 \(r\) 的嵌入
\[
j:\ZZ^r\hookrightarrow A.
\]
由于 \(A\cong\QQ^r\) 且 \(j(\ZZ^r)\) 的秩已等于 \(r\)，其 \(\QQ\)-线性张成即为整个 \(A\)。
于是局部化的泛性质把 \(j\) 唯一延拓为一个同构
\[
\phi:\QQ^r\overset{\sim}{\longrightarrow}A.
\]
对偶化得到唯一连续同构
\[
\widehat{\phi}:K\overset{\sim}{\longrightarrow}\widehat{\QQ}^{\,r}
\]
并满足 \(\pi_{\QQ,r}\circ \widehat{\phi}=\pi\)。
故 \((\widehat{\QQ}^{\,r},\pi_{\QQ,r})\) 为所述范畴的终对象。
证毕。
\end{proof}

\begin{corollary}[通用连通包络相对于基底的自同构刚性]\label{cor:conclusion-boundary-solenoid-over-base-rigidity}
设
\[
u:\widehat{\QQ}^{\,r}\longrightarrow \widehat{\QQ}^{\,r}
\]
是连续群同态，并满足
\[
\pi_{\QQ,r}\circ u=\pi_{\QQ,r},
\]
其中
\[
0\to \widehat{\ZZ}^{\,r}\to \widehat{\QQ}^{\,r}\xrightarrow{\ \pi_{\QQ,r}\ }\TT^r\to 0
\]
为定理 \ref{thm:conclusion-boundary-adelic-solenoid-envelope} 的典范短正合列。则必有
\[
u=\mathrm{id}_{\widehat{\QQ}^{\,r}}.
\]
\end{corollary}
\begin{proof}
由定理 \ref{thm:conclusion-boundary-adelic-solenoid-envelope}，\((\widehat{\QQ}^{\,r},\pi_{\QQ,r})\) 在相应范畴中是终对象。映射 \(u\) 与恒等映射
\(
\mathrm{id}_{\widehat{\QQ}^{\,r}}
\)
都是从该终对象到自身、并与到 \(\TT^r\) 的投影相容的态射。终对象上的此类态射唯一，故 \(u=\mathrm{id}_{\widehat{\QQ}^{\,r}}\)。证毕。
\end{proof}

\begin{theorem}[阿德尔边界 torsor 的有限素数角色分解]\label{thm:conclusion-boundary-adelic-torsor-finite-prime-character-decomposition}
在定理 \ref{thm:conclusion-boundary-adelic-solenoid-envelope} 的设定下，选定 \(\mathcal F_{\widehat{\ZZ}}(\delta)\) 的一个基点，并识别
\[
\mathcal F_{\widehat{\ZZ}}(\delta)\cong \widehat{\ZZ}^{\,r}=\prod_p \ZZ_p^{\,r}.
\]
则其 Pontryagin 对偶满足
\[
\widehat{\mathcal F_{\widehat{\ZZ}}(\delta)}
\cong
\bigoplus_p C_{p^\infty}^{\,r},
\]
其中 \(C_{p^\infty}\) 为 Prüfer \(p\)-群。特别地，每个连续角色 \(\chi\) 都唯一分解为有限和
\[
\chi=\sum_{p\in S_\chi}\chi_p,
\qquad
S_\chi\subset \mathbb P\ \text{有限},
\qquad
\chi_p\in C_{p^\infty}^{\,r}.
\]
因此，阿德尔边界 torsor 上不存在同时看见无穷多个素数方向的单一连续角色。
\end{theorem}
\begin{proof}
Pontryagin 对偶把紧直积送到离散直和，故
\[
\widehat{\mathcal F_{\widehat{\ZZ}}(\delta)}
\cong
\widehat{\widehat{\ZZ}^{\,r}}
\cong
\widehat{\prod_p \ZZ_p^{\,r}}
\cong
\bigoplus_p \widehat{\ZZ_p^{\,r}}.
\]
又 \(\widehat{\ZZ_p}\cong C_{p^\infty}\)，从而
\[
\widehat{\ZZ_p^{\,r}}\cong C_{p^\infty}^{\,r}.
\]
故得到
\[
\widehat{\mathcal F_{\widehat{\ZZ}}(\delta)}
\cong
\bigoplus_p C_{p^\infty}^{\,r}.
\]
直和中的元素只有有限多个非零分量，因此每个连续角色都具有唯一的有限素数支撑分解。证毕。
\end{proof}

\begin{theorem}[有限维连续相位通道对无穷素数尾的结构性失明]\label{thm:conclusion-boundary-prime-tail-blindness}
任取连续群同态
\[
L:\widehat{\ZZ}^{\,r}\longrightarrow \TT^d.
\]
则存在有限素数集 \(S_L\subset \mathbb P\)，使得 \(L\) 因子化为
\[
\widehat{\ZZ}^{\,r}\longrightarrow \prod_{p\in S_L}\ZZ_p^{\,r}\xrightarrow{\ \widetilde L\ }\TT^d.
\]
等价地，
\[
\prod_{p\notin S_L}\ZZ_p^{\,r}\subseteq \ker L.
\]
因此，任何有限维连续相位通道都只能看见有限多个素数方向，而无穷素数尾必然落入其不可见核中。
\end{theorem}
\begin{proof}
对偶化得到离散群同态
\[
\widehat L:\ZZ^d\longrightarrow \widehat{\widehat{\ZZ}^{\,r}}
\cong
\bigoplus_p C_{p^\infty}^{\,r}.
\]
由于 \(\ZZ^d\) 有限生成，其像 \(\widehat L(\ZZ^d)\) 亦有限生成。直和群
\(
\bigoplus_p C_{p^\infty}^{\,r}
\)
中的任意有限生成子群只能支撑在有限多个素数分量上，因为有限组生成元各自只有有限支撑，而有限个有限集之并仍有限。故存在有限素数集 \(S_L\) 使
\[
\widehat L(\ZZ^d)\subseteq \bigoplus_{p\in S_L}C_{p^\infty}^{\,r}.
\]
再对偶回去，即得 \(L\) 仅通过 \(\prod_{p\in S_L}\ZZ_p^{\,r}\) 因子化。于是全部 \(p\notin S_L\) 的方向都落入 \(\ker L\)。证毕。
\end{proof}

\leanverified{paper\_conclusion\_boundary\_minimal\_phase\_circle\_count}
\begin{theorem}[边界解纤维的最小相位圆数]\label{thm:conclusion-boundary-minimal-phase-circle-count}
固定素数 \(p\) 与层级 \(\ell\ge 1\)，并取一条非空纤维 \(\mathcal F_{p^\ell}(\delta_\ell)\)。
选定基点与 \((\ZZ/p^\ell\ZZ)\)-基后，将其识别为
\[
Q_{p,\ell}^{(r)}:=\{0,1,\dots,p^\ell-1\}^{r}.
\]
设
\[
\Phi_\ell:Q_{p,\ell}^{(r)}\longrightarrow \TT^d
\]
是一种纯相位读出，并以 \(\TT^d\) 的乘积 sup 距离 \(d_\infty\) 衡量精度。
若存在 \(\delta_\ell\in(0,\tfrac14)\) 使得 \(\Phi_\ell(Q_{p,\ell}^{(r)})\) 中任意两点皆满足
\[
d_\infty\bigl(\Phi_\ell(u),\Phi_\ell(v)\bigr)>2\delta_\ell
\qquad(u\neq v),
\]
则必存在仅依赖于 \(d\) 的常数 \(C_d>0\) 使
\[
\delta_\ell\le C_d\,p^{-\ell r/d}.
\]
等价地，令
\[
b_\ell:=\Bigl\lceil \log_2\frac{1}{\delta_\ell}\Bigr\rceil,
\]
则
\[
\boxed{\;
b_\ell\ \ge\ \frac{r}{d}\,\ell\log_2 p-O(1).
\;}
\]
特别地，若要求相位读出达到原生 \(p^\ell\)-分辨率，即
\[
b_\ell=\ell\log_2 p+O(1)
\qquad\text{（等价地 }\delta_\ell\asymp p^{-\ell}\text{）},
\]
则必要且充分条件为
\[
d\ge r.
\]
从而边界方程解纤维的最小相位圆数恰为
\[
\boxed{\;d_{\min}=r(f).\;}
\]
\end{theorem}
\begin{proof}
由 \(d_\infty\)-距离下的体积打包估计，\(\TT^d\) 中任意两两 \(2\delta_\ell\)-分离的点集基数至多为 \(C_d'\delta_\ell^{-d}\)（常数仅依赖于 \(d\)）。
而此处点集大小为
\[
\card{Q_{p,\ell}^{(r)}}=p^{\ell r},
\]
故必须有
\[
p^{\ell r}\le C_d'\delta_\ell^{-d},
\]
即
\[
\delta_\ell\le C_d\,p^{-\ell r/d}
\]
（吸收常数后得到）。
再由 \(b_\ell=\lceil \log_2(1/\delta_\ell)\rceil\) 立即推出位深下界
\[
b_\ell\ge \frac{r}{d}\,\ell\log_2 p-O(1).
\]
这也与定理 \ref{thm:cdim-prefix-entropy-depth-conservation} 在 \(G=\ZZ^r\)、\(N=p^\ell-1\) 的专门化一致。

若 \(b_\ell=\ell\log_2 p+O(1)\)，则上式强制 \(d\ge r\)。
反之当 \(d=r\) 时，标准嵌入
\[
\iota_{p^\ell}:Q_{p,\ell}^{(r)}\hookrightarrow \TT^r,
\qquad
(a_1,\dots,a_r)\longmapsto
\left(\frac{a_1}{p^\ell},\dots,\frac{a_r}{p^\ell}\right)\bmod \ZZ^r
\]
的像点最小 \(d_\infty\)-间距恰为 \(p^{-\ell}\)。
因此当 \(\delta_\ell<\tfrac12 p^{-\ell}\) 时可唯一恢复，说明 \(d=r\) 可达。
故最小相位圆数为 \(r\)。
证毕。
\end{proof}

\begin{theorem}[有限素数联合读出的纯相位打包律]\label{thm:conclusion-boundary-finite-prime-pure-phase-packing}
设 \(S\subset \mathbb P\) 为有限素数集，\(\boldsymbol\ell=(\ell_p)_{p\in S}\) 为正整数向量，并记
\[
M_S:=\prod_{p\in S}p^{\ell_p}.
\]
固定一条非空联合纤维 \(\mathcal F_{S,\boldsymbol\ell}(\delta)\)，选定基点并经中国剩余定理把它识别为
\[
Q_{M_S}^{(r)}:=\{0,1,\dots,M_S-1\}^r\cong (\ZZ/M_S\ZZ)^r.
\]
设
\[
\Phi_{S,\boldsymbol\ell}:Q_{M_S}^{(r)}\longrightarrow \TT^d
\]
是一条纯相位读出，并以 \(\TT^d\) 的乘积 sup 距离 \(d_\infty\) 衡量精度。若存在 \(\delta\in(0,\tfrac14)\) 使得像点两两满足
\[
d_\infty\bigl(\Phi_{S,\boldsymbol\ell}(u),\Phi_{S,\boldsymbol\ell}(v)\bigr)>2\delta
\qquad(u\neq v),
\]
则存在仅依赖于 \(d\) 的常数 \(C_d>0\) 使
\[
\delta\le C_d\,M_S^{-r/d}.
\]
等价地，令
\[
b:=\Bigl\lceil \log_2\frac{1}{\delta}\Bigr\rceil,
\]
则
\[
\boxed{\;
d\,b\ge r\log_2 M_S-O(1).
\;}
\]
特别地，在原生分辨率
\[
b=\log_2 M_S+O(1)
\]
下，必要且充分条件为
\[
d\ge r.
\]
\end{theorem}
\begin{proof}
由 \(d_\infty\)-距离下的体积打包估计，\(\TT^d\) 中任意两两 \(2\delta\)-分离点集的基数至多为 \(C_d'\delta^{-d}\)。
而
\[
\card{Q_{M_S}^{(r)}}=M_S^r,
\]
故必须有
\[
M_S^r\le C_d'\delta^{-d},
\]
整理即得
\[
\delta\le C_d\,M_S^{-r/d}.
\]
取二进对数后便得到
\[
d\,b\ge r\log_2 M_S-O(1).
\]

若 \(b=\log_2 M_S+O(1)\)，则上式强制 \(d\ge r\)。
反之当 \(d=r\) 时，标准嵌入
\[
Q_{M_S}^{(r)}\hookrightarrow \TT^r,
\qquad
(a_1,\dots,a_r)\longmapsto
\left(\frac{a_1}{M_S},\dots,\frac{a_r}{M_S}\right)\bmod \ZZ^r
\]
的像点最小 \(d_\infty\)-间距恰为 \(M_S^{-1}\)。因此当 \(\delta<\tfrac12 M_S^{-1}\) 时可唯一恢复，说明 \(d=r\) 可达。证毕。
\end{proof}

\begin{theorem}[离散寄存器与相位圆的混合预算律]\label{thm:conclusion-boundary-mixed-register-phase-budget}
\leanverified{paper\_conclusion\_boundary\_mixed\_register\_phase\_budget}
沿用定理 \ref{thm:conclusion-boundary-minimal-phase-circle-count} 的设定。
设 \(R_\ell\) 为有限离散寄存器，满足
\[
\card{R_\ell}\le 2^{B_\ell},
\]
并给定编码
\[
\mathrm{Enc}_\ell:Q_{p,\ell}^{(r)}\longrightarrow R_\ell\times \TT^d.
\]
若存在 \(\delta_\ell\in(0,\tfrac14)\) 使得对每个固定的寄存器值 \(\rho\in R_\ell\)，相位子集
\[
\left\{\theta\in\TT^d:\ (\rho,\theta)\in \mathrm{Enc}_\ell\bigl(Q_{p,\ell}^{(r)}\bigr)\right\}
\]
中的不同点两两 \(2\delta_\ell\)-分离，则记
\[
b_\ell:=\Bigl\lceil \log_2\frac{1}{\delta_\ell}\Bigr\rceil,
\]
必有线性预算下界
\[
\boxed{\;
B_\ell+d\,b_\ell\ \ge\ r\,\ell\log_2 p-O(1).
\;}
\]
进一步，若离散寄存器由 \(a\) 条 \(p\)-进轴、每条深度 \(\ell\) 实现，即
\[
\card{R_\ell}=p^{a\ell+O(1)}
\qquad\text{（等价地 }B_\ell=a\,\ell\log_2 p+O(1)\text{）},
\]
则
\[
\boxed{\;
d\,b_\ell\ \ge\ (r-a)\,\ell\log_2 p-O(1).
\;}
\]
在原生相位分辨率 \(b_\ell=\ell\log_2 p+O(1)\) 下，上式化为
\[
\boxed{\;a+d\ \ge\ r.\;}
\]
并且该下界是可达的：对任意 \(0\le a\le r\)，取前 \(a\) 个坐标直接写入 \(a\) 条 \(p\)-进寄存轴，余下 \(r-a\) 个坐标按标准嵌入写入 \(\TT^{r-a}\)，即可在 \(a+d=r\) 时实现原生分辨率的唯一恢复。
\end{theorem}
\begin{proof}
由鸽巢原理，存在某个寄存器值 \(\rho_\ast\in R_\ell\)，使得其对应的相位纤维至少含有
\[
\frac{p^{\ell r}}{\card{R_\ell}}\ \ge\ p^{\ell r}\,2^{-B_\ell}
\]
个点。
这些点在 \(\TT^d\) 中两两 \(2\delta_\ell\)-分离。
再次应用 \(\TT^d\) 的体积打包估计，存在常数 \(C_d'>0\) 使
\[
p^{\ell r}\,2^{-B_\ell}\le C_d'\delta_\ell^{-d}.
\]
取 \(\log_2\) 并代入 \(b_\ell=\lceil \log_2(1/\delta_\ell)\rceil\)，得
\[
r\,\ell\log_2 p-B_\ell\le d\,b_\ell+O(1),
\]
即第一条预算下界。

若 \(\card{R_\ell}=p^{a\ell+O(1)}\)，则
\(
B_\ell=a\ell\log_2 p+O(1)
\)，代回即得
\[
d\,b_\ell\ge (r-a)\ell\log_2 p-O(1).
\]
进一步取 \(b_\ell=\ell\log_2 p+O(1)\) 即推出 \(a+d\ge r\)。

反向可达性由显式分裂构造给出：把
\[
Q_{p,\ell}^{(r)}=
\{0,\dots,p^\ell-1\}^{a}\times \{0,\dots,p^\ell-1\}^{r-a}
\]
的前 \(a\) 个坐标保留为 \(p\)-进寄存器，后 \(r-a\) 个坐标用
\[
(b_1,\dots,b_{r-a})\longmapsto
\left(\frac{b_1}{p^\ell},\dots,\frac{b_{r-a}}{p^\ell}\right)\bmod \ZZ^{r-a}
\]
写入 \(\TT^{r-a}\)。
此时相位最小间距为 \(p^{-\ell}\)，故在 \(\delta_\ell<\tfrac12p^{-\ell}\) 时唯一可逆，并且恰有 \(a+d=r\)。
证毕。
\end{proof}

\begin{theorem}[有限素数联合读出的 sharp mixed budget frontier]\label{thm:conclusion-boundary-finite-prime-mixed-pareto-frontier}
\leanverified{paper\_conclusion\_boundary\_finite\_prime\_mixed\_pareto\_frontier}
设 \(S\subset \mathbb P\) 为有限素数集，\(\boldsymbol\ell=(\ell_p)_{p\in S}\) 为正整数向量，并记
\[
M_S:=\prod_{p\in S}p^{\ell_p}.
\]
固定一条非空联合纤维并经中国剩余定理识别为
\[
Q_{M_S}^{(r)}\cong (\ZZ/M_S\ZZ)^r.
\]
设 \(R\) 为有限离散寄存器，满足 \(\card{R}\le 2^B\)，并给定编码
\[
\mathrm{Enc}_{S,\boldsymbol\ell}:Q_{M_S}^{(r)}\longrightarrow R\times \TT^d.
\]
若存在 \(\delta\in(0,\tfrac14)\) 使得对每个固定寄存器值 \(\rho\in R\)，对应相位子集中的不同点两两 \(2\delta\)-分离，并记
\[
b:=\Bigl\lceil \log_2\frac{1}{\delta}\Bigr\rceil,
\]
则必有
\[
\boxed{\;
B+d\,b\ge r\log_2 M_S-O(1).
\;}
\]
进一步，若某个整数 \(a\in\{0,1,\dots,r\}\) 满足
\[
B=a\log_2 M_S+O(1),
\]
且要求原生相位分辨率
\[
b=\log_2 M_S+O(1),
\]
则必要且充分条件为
\[
a+d\ge r.
\]
并且极小可达前沿恰为
\[
\boxed{\;a+d=r.\;}
\]
\end{theorem}
\begin{proof}
由鸽巢原理，存在某个寄存器值 \(\rho_\ast\in R\)，使得对应相位纤维至少含有
\[
\frac{M_S^r}{\card{R}}
\ge
M_S^r\,2^{-B}
\]
个点。对该纤维应用定理 \ref{thm:conclusion-boundary-finite-prime-pure-phase-packing} 的打包估计，得
\[
M_S^r\,2^{-B}\le C_d'\delta^{-d}.
\]
取二进对数即得
\[
B+d\,b\ge r\log_2 M_S-O(1).
\]

若 \(B=a\log_2 M_S+O(1)\) 且 \(b=\log_2 M_S+O(1)\)，则上式化为
\[
a+d\ge r.
\]
反向可达性由显式坐标分裂给出：把
\[
Q_{M_S}^{(r)}=\{0,\dots,M_S-1\}^{a}\times \{0,\dots,M_S-1\}^{r-a}
\]
的前 \(a\) 个坐标直接写入离散寄存器，后 \(r-a\) 个坐标按
\[
(b_1,\dots,b_{r-a})\longmapsto
\left(\frac{b_1}{M_S},\dots,\frac{b_{r-a}}{M_S}\right)\bmod \ZZ^{r-a}
\]
写入 \(\TT^{r-a}\)。
此时寄存器规模为 \(M_S^a\)，相位最小间距为 \(M_S^{-1}\)，故在 \(\delta<\tfrac12 M_S^{-1}\) 时可唯一恢复。于是当 \(d=r-a\) 时恰达到原生分辨率，故极小前沿为 \(a+d=r\)。证毕。
\end{proof}

\begin{theorem}[面积密度非退化族上的 native budget 爆炸律]\label{thm:conclusion-boundary-native-budget-explosion}
取一列满射粗粒化 \(f_n:\Omega_n\twoheadrightarrow X_n\)，并记面积密度
\[
a_n:=\frac{A(f_n)}{n2^{n-1}}.
\]
若
\[
\liminf_{n\to\infty}a_n>0,
\qquad
\frac{\card{X_n}}{2^n}\to 0,
\]
则由推论 \ref{cor:conclusion-boundary-cycle-rank-density-explosion}，
\[
r_n:=r(f_n)=\Theta(n2^{n-1}).
\]
固定任意素数 \(p\)。若在第 \(n\) 层、精度 \(\ell_n\) 处采用编码
\[
\mathrm{Enc}_{n,\ell_n}:(\ZZ/p^{\ell_n}\ZZ)^{r_n}\longrightarrow R_n\times \TT^{d_n},
\]
并要求原生相位分辨率
\[
b_n=\ell_n\log_2 p+O(1),
\]
同时存在整数 \(\alpha_n^{\mathrm{reg}}\ge 0\) 满足
\[
\card{R_n}=p^{\alpha_n^{\mathrm{reg}}\ell_n+O(1)},
\]
则必有
\[
\boxed{\;
\alpha_n^{\mathrm{reg}}+d_n\ge r_n.
\;}
\]
特别地：
\begin{enumerate}
  \item 若 \(d_n\le D\) 一致有界，则
  \[
  \log_2\card{R_n}
  \ge
  (r_n-D)\ell_n\log_2 p-O(1)
  =
  \Omega(\ell_n n2^{n-1});
  \]
  \item 若 \(\alpha_n^{\mathrm{reg}}\le A\) 一致有界，则
  \[
  d_n\ge r_n-A=\Theta(n2^{n-1}).
  \]
\end{enumerate}
\end{theorem}
\begin{proof}
由推论 \ref{cor:conclusion-boundary-cycle-rank-density-explosion}，
\(
r_n=\Theta(n2^{n-1})
\)。
对每个固定的 \(n\)，把定理 \ref{thm:conclusion-boundary-finite-prime-mixed-pareto-frontier} 专门化到单素数集合 \(S=\{p\}\) 与 \(M_S=p^{\ell_n}\)。
此时
\[
B_n=\log_2\card{R_n}
=
\alpha_n^{\mathrm{reg}}\ell_n\log_2 p+O(1),
\qquad
b_n=\ell_n\log_2 p+O(1),
\]
故由该定理的原生前沿结论立即得到
\[
\alpha_n^{\mathrm{reg}}+d_n\ge r_n.
\]

若 \(d_n\le D\)，则
\[
\alpha_n^{\mathrm{reg}}\ge r_n-D,
\]
从而
\[
\log_2\card{R_n}
=
\alpha_n^{\mathrm{reg}}\ell_n\log_2 p+O(1)
\ge
(r_n-D)\ell_n\log_2 p-O(1)
=
\Omega(\ell_n n2^{n-1}).
\]
若 \(\alpha_n^{\mathrm{reg}}\le A\)，则直接得到
\[
d_n\ge r_n-A=\Theta(n2^{n-1}).
\]
证毕。
\end{proof}

\begin{conjecture}[native 最优编码的 CRT 分裂刚性]\label{conj:conclusion-boundary-native-optimal-encoder-crt-rigidity}
固定有限素数集 \(S\subset\mathbb P\)、精度向量 \(\boldsymbol\ell\) 与秩 \(r\)，并记
\[
M_S:=\prod_{p\in S}p^{\ell_p}.
\]
设
\[
\mathrm{Enc}_{S,\boldsymbol\ell}:(\ZZ/M_S\ZZ)^r\longrightarrow R\times \TT^d
\]
是一条 exact native encoder，并达到定理 \ref{thm:conclusion-boundary-finite-prime-mixed-pareto-frontier} 的极小前沿，即
\[
a+d=r,
\qquad
\card{R}=M_S^a.
\]
猜想：在下列等价关系下，\(\mathrm{Enc}_{S,\boldsymbol\ell}\) 与标准 CRT 坐标分裂模型唯一等价：
\begin{enumerate}
  \item 源空间允许 \(\operatorname{GL}_r(\ZZ/M_S\ZZ)\) 线性自同构；
  \item 寄存器允许有限集合双射；
  \item 相位端允许由 \(\operatorname{GL}_d(\ZZ)\) 诱导的 torus 自同构与平移。
\end{enumerate}
亦即，每一条 saturating encoder 都应当等价于
\[
(x_1,\dots,x_r)\longmapsto
\Bigl(
(x_1,\dots,x_a),\ 
(x_{a+1}/M_S,\dots,x_r/M_S)\bmod \ZZ^d
\Bigr).
\]
换言之，不存在 genuinely nonlinear 的最优替代机制。
\end{conjecture}

\paragraph{结构闭合}
就边界方程而言，回路秩同时控制有限层解纤维、逐层精确选择器、有限素数联合预算、其 \(p\)-进与阿德尔完成、通用连通包络、有限素数支撑的角色分解以及 native mixed architecture 的爆炸尺度。
因此在统一口径下得到结构链
\[
\mathsf B_f
\Longrightarrow
r(f)
\Longrightarrow
\mathcal F_{p^\ell}(\delta_\ell)\cong (\ZZ/p^\ell\ZZ)^{r(f)}
\Longrightarrow
\card{\mathcal F_{p^\ell}(\delta_\ell)}=p^{\ell r(f)}
\Longrightarrow
\min\card{R_\ell}=p^{\ell r(f)}
\Longrightarrow
\mathcal F_{S,\boldsymbol\ell}(\delta)\cong (\ZZ/M_S\ZZ)^{r(f)}
\]
\[
\Longrightarrow
\widehat{\ZZ}^{\,r(f)}\text{-torsor}
\Longrightarrow
\widehat{\QQ}^{\,r(f)}\xrightarrow{\ \pi_{\QQ,r(f)}\ }\TT^{\,r(f)}
\Longrightarrow
\widehat{\mathcal F_{\widehat{\ZZ}}(\delta)}
\cong
\bigoplus_p C_{p^\infty}^{\,r(f)}
\Longrightarrow
\text{有限维相位通道仅见有限素数支撑}
\Longrightarrow
d_{\min}=r(f)
\Longrightarrow
B+d\,b\ge r(f)\log_2 M_S-O(1)
\Longrightarrow
\text{原生前沿 }a+d=r(f)
\Longrightarrow
\text{面积密度非退化时 native budget 爆炸}.
\]

\endinput


\endinput
